\documentclass{beamer}

\usetheme{UCL2023}

%%%%%%%%%%%%%%%%%%%%%%%%%%%%%%%%%%%%%%%%%%%%%%%%%%%%%
%%%%%%%%%%%%%%%%%%%%%%%%%%%%%%%%%%%%%%%%%%%%%%%%%%%%%%%%             

\usepackage{graphicx}
\usepackage{amssymb}
\usepackage{amsmath}

\graphicspath{ {./images/} }

%%%%%%%%%%%%%%%%%%%%%%%%%%%%%%%%%%%%%%%%%%%%%%%%%%%%%%


\title{Multiple imputation: brief summary}
%\subtitle{with Stata and R}
\author{Aur\'elie Bertrand \& Laura Symul}
\date{Applied Statistics Workshop, Nov 2026}
\institute{SMCS, ISBA} 
\begin{document}


%%%%%%%%%%%%%%%%%%%%%%%%%%%%%%%%%%%%%%%%%%%%%%%%%%%%%%

\begin{frame}[plain]
 \titlepage
\end{frame}

%%%%%%%%%%%%%%%%%%%%%%%%%%%%%%%%%%%%%%%%%%%%%%%%%%%%%%

\section{Multiple imputation (MI)}


\begin{frame}
\frametitle{A 3-step method}

Each missing cell is \textcolor{cyan}{replaced by more than one value}. %, the analysis is performed on each imputed dataset and the results are then \textcolor{cyan}{pooled} over all these datasets.

\vspace{0.3cm}

		\begin{enumerate}
			\item \textcolor{cyan}{Imputation} step 
				\begin{itemize}
					\item Create $m>1$ imputed (and, hence, complete) datasets by replacing, for each variable, all missing values by plausible values from a specifically designed distribution.
				\end{itemize}
			\item \textcolor{cyan}{Analysis} (or estimation) step
				\begin{itemize}
					\item Apply the traditional (complete-data) analysis (\textit{e.g.}, linear regression) on each imputed dataset.
				\end{itemize}
			\item \textcolor{cyan}{Pooling} step
				\begin{itemize}
					\item Results from the different data sets are pooled to obtain the final estimates. 
				\end{itemize}
		\end{enumerate}
\end{frame}

\begin{frame}
\frametitle{Advantages}

	\begin{itemize}
	\item Requires \textcolor{cyan}{MAR}.% (or MCAR).	
	\item Takes into account the \textcolor{cyan}{uncertainty in the imputed values}.
			\begin{itemize}
				\item Contrary to single imputation.
			\end{itemize}
	\item Incorporates \textcolor{cyan}{auxiliary information} (in the imputation step).

	\item With plausible imputation model, produces \textcolor{cyan}{unbiased parameter and standard error} estimators. 
		\begin{itemize}
			\item Quite \textcolor{cyan}{robust} to departures from the correct specification of the imputation model if the rates of missing information are low and the correct model of interet is used in the analysis.
		\end{itemize}			
	\item Allows for \textcolor{cyan}{varying uncertainty levels} between imputed variables.
	\item The missing value problem is dealt with \textcolor{cyan}{before the analysis} of interest is performed.
	\item Can be used with many types of analyses.
\end{itemize}
	
	
\end{frame}


\begin{frame}
\frametitle{MI in practice}
In `R`
\begin{itemize}
	\item In general, performed using three separate steps:
		\begin{itemize}
			\item Generate the imputed datasets (using, for example, the \texttt{mice} function of package \texttt{mice} or the \texttt{amelia} function of package \texttt{Amelia}).
			\item Analyze them (using, for example, the \texttt{with} function of the \texttt{mice} package).
			\item Pool the results (using, for example, the \texttt{pool} function of the \texttt{mice} package).
		\end{itemize}
\end{itemize}

Many alternative implementations exist for different software/language.

\end{frame}


\subsection{Imputation step}

\begin{frame}
\frametitle{Imputation step -- Main parameters}
\begin{itemize}
	\item \textcolor{cyan}{Number of imputations}, $m$:
		\begin{itemize}
			\item (rule of thumb) set $m$ to the average (over all variables) percentage of missing data.
			\item Recommended to use at least 20 imputations to reduce sampling error.
		\end{itemize}
	\item Choice of the \textcolor{cyan}{imputation method}.
		\begin{itemize}
			\item If only 1 variable with missing values: choose a \textcolor{cyan}{univariate} imputation method and its \textcolor{cyan}{predictors}. 
			\item If missing values in more than one variable: choose between
				\begin{itemize}
					\item an approach which imputes \textcolor{cyan}{all variables at the same time}.
					\item an approach that imputes the variables \textcolor{cyan}{one by one} (iteratively): for each variable, a univariate imputation method must be chosen.
				\end{itemize}
	
		\end{itemize}
\end{itemize}
\vfill
\footnotesize{\url{https://stefvanbuuren.name/fimd/sec-howmany.html}}
\end{frame}


\subsubsection{Univariate methods}


\begin{frame}
\frametitle{Univariate imputation methods}
A \textcolor{cyan}{statistical model} allowing one to \textcolor{cyan}{predict the variable with missing values} based on other variables from the data.\\ \vspace{0.3cm}

In multiple imputation, it is therefore important to distinguish between:

\begin{itemize}
	\item the \textcolor{cyan}{response model}: the model of interest.
	\item the \textcolor{cyan}{imputation model}: the model used to impute the incomplete variable.
\end{itemize}


\end{frame}



\begin{frame}
\frametitle{Choice of a univariate imputation method}

Choice based on the \textcolor{cyan}{type of the variable} to impute:
\begin{itemize}
	\item Continuous variable:
		\begin{itemize}
			\item \textcolor{cyan}{Bayesian multiple imputation}
			\item \textcolor{cyan}{Bootstrap multiple imputation}
			\item \textcolor{cyan}{Predictive mean matching} (PMM): less sensitive to model misspecification (\textit{e.g.}, nonnormality) than Bayesian and Bootstrap multiple imputation.
	\end{itemize}
	\item Discrete or categorical variable:
		\begin{itemize}
			\item Multiple imputation based on binary, multinomial or ordered \textcolor{cyan}{logistic regression}.
			\item \textcolor{cyan}{Predictive mean matching} (PMM)
		\end{itemize}
	\item Count variables, semi-continuous variables, censored variables, truncated variables, rounded variables: some methods exist too.
\end{itemize}

\emph{An in-depth comparison of these methods (useful for making an informed choice) is beyond the scope of this workshop.}

\end{frame}


\begin{frame}
\frametitle{Predictive Mean Matching (PMM)}		

	\begin{itemize}
		\item Alternative to Bayesian and Bootstrap multiple imputation.
		\item Preserves the original distribution $\Rightarrow$ \textcolor{cyan}{less sensitive to model misspecification} (\textit{e.g.}, nonnormality) than Bayesian and Bootstrap multiple imputation.

		\item Quite robust to transformations of the incomplete variable (generally yields similar results, whatever the transformation).
		\item Suited to \textcolor{cyan}{discrete data}.
	
	\end{itemize}
	\vfill
{\footnotesize Jia and Wu (2023), \url{https://stefvanbuuren.name/fimd/sec-pmm.html}}
\end{frame}

\begin{frame}[label=main3]
\frametitle{Predictive Mean Matching (PMM) (cont.)}		

Steps:
	\begin{itemize}
		\item For each variable containing missing values: 
			\begin{itemize}
				\item Estimate the parameters of a \textcolor{cyan}{predictive model}.
				\item For each (missing or observed) value, \textcolor{cyan}{obtain a prediction} from this estimated predictive model.
				\item Then, for each missing value:
					\begin{itemize}
						\item Identify a set of \textcolor{cyan}{candidate donors}, i.e., a set of observed values whose predicted value is close to the predicted value of the missing value. 
						\item Randomly \textcolor{cyan}{select one donor} among the candidate donors, and use its observed value as \textcolor{cyan}{imputed value}.
					\end{itemize}
			\end{itemize}
	\end{itemize}
	In R, \texttt{mice} function with the \texttt{pmm} method. By default, a linear model with main effects of all other variables.


%\hyperlink{supplemental3}{\beamerbutton{More details}}\\

\vfill
\footnotesize{Jia and Wu (2023), \url{https://stefvanbuuren.name/fimd/sec-pmm.html}}


\end{frame}


\begin{frame}
%\frametitle{Imputation methods for univariate missing data}
\frametitle{Univariate imputation methods in common software}
In R: many methods available, among which:
\begin{itemize}
	\item Bayesian multiple imputation: in \texttt{mice}, method \texttt{norm} (or, equivalently, \texttt{mice.impute.norm}).
	\item Bootstrap multiple imputation: in \texttt{mice}, method \texttt{norm.boot}.
	\item Predictive mean matching: in \texttt{mice}, method \texttt{pmm} (default method for continuous data).
	\item Logistic regressions: methods \texttt{logreg} (for binary variables), \texttt{polyreg} (for unordered variables) and \texttt{polr} (for ordered variables).
\end{itemize}

\end{frame}



\subsection{Predictor variables in the imputation model}

\begin{frame}
\frametitle{Predictor variables in the imputation model}

The imputation model should contain:
\begin{itemize}
\setlength\itemsep{0.3cm}
	\item All variables involved in the definition of the \textcolor{cyan}{model of interest} and the \textcolor{cyan}{sampling design}:
		\begin{itemize}
			\item All variables that are \textcolor{cyan}{explanatory variables} in the model of interest (including interactions).
			\item The \textcolor{cyan}{response variable} of the model of interest.	
			\item In a survey context: sampling weights, (main) strata and (main) cluster identifiers.
		\end{itemize}
		 \emph{Congeniality} principle: all variables of the analysis model should also appear (in the same form) in the imputation model.
	\item \textcolor{cyan}{Auxiliary variables}:
		\begin{itemize}
			\item All variables that are related to missingness.
			\item All variables that are related to the variable with missing values.
		\end{itemize}

\end{itemize}

\vfill
\footnotesize{Mainzer et al. (2021), Stata (2023a)}
\end{frame}


\begin{frame}
\frametitle{Choice of predictor variables in the imputation model}
\begin{itemize}
\setlength\itemsep{0.3cm}
	\item If max 20-30 variables: include all of them.
	\item If more variables: select 15-20 variables:
		\begin{itemize}
			\item All the variables of the \textcolor{cyan}{analysis model} (including the outcome and interactions).
			\item \textcolor{cyan}{Auxiliary} variables (related to nonresponse and to the variable to be imputed).
			\item \textcolor{cyan}{Omit variables having too many missing values} in the subgroup of observations where the variable to be imputed is missing.
		\end{itemize}
	\item There exist \textcolor{cyan}{indices} that can be computed based on the data and that can help selecting predictor variables (\emph{beyond the scope of this introductory training}).
\end{itemize}
\vfill
\footnotesize{\url{https://stefvanbuuren.name/fimd/sec-modelform.html}}

\end{frame}


\subsection{Pooling phase}



\begin{frame}
\frametitle{MI: pooling phase for parameter estimates}
Combine estimated parameters and their SEs using rules that allow overall inference (Rubin, 1987), known as \textcolor{cyan}{Rubin's rules}.
\begin{itemize}
\setlength\itemsep{0.3cm}
	\item Final estimation of a \textcolor{cyan}{parameter}: 
		\begin{itemize}
			\item average (over the imputed datasets) of all estimated values.
		\end{itemize}
	\item Final estimation of a parameter \textcolor{cyan}{standard error}: 
		\begin{itemize}
			\item NOT the average of all standard errors!
			\item Average of two sources of variability:
		
			\begin{itemize}
				\item Average of \textcolor{cyan}{within-imputation variance}: average (over the imputed datasets) of variances in the estimation of a parameter in a sample (related to its SE).
				\item \textcolor{cyan}{Between-imputation variance}: variability (between imputed datasets) in the estimations of the parameter.
			\end{itemize}
		\end{itemize}
\end{itemize}
\vfill
\footnotesize{\url{https://stefvanbuuren.name/fimd/sec-multiparameter.html}}
\end{frame}


\begin{frame}
\frametitle{MI: pooling phase for tests}

Confidence intervals and p-values of a test about a \textcolor{cyan}{single parameter}:
\begin{itemize}
	\item NOT the average of all confidence interval bounds or p-values!
	\item Computed based on pooled estimation of the parameter and its standard error (as seen on previous slide).
\end{itemize}	

\vspace{1cm}

For \textcolor{cyan}{multiparameter} tests, \textit{e.g.} test of several coefficients, test of model fit, test of comparison of fit of two models:  
\begin{itemize}
	\item 3 possible types of test statistics.
\end{itemize}	
\vfill
\footnotesize{\url{https://stefvanbuuren.name/fimd/sec-multiparameter.html}}
\end{frame}



\subsection{Miscellaneous}



\begin{frame}
\frametitle{MI in practice}

\begin{enumerate}
	\item Decide whether MAR is plausible
	\item Choose the form of the structural part of the imputation model and the error distribution.
	\item Select variables to include in the imputation model. %General advice: as many variables (and their interactions!) as possible.
	\item Decide whether to impute variables that are functions of other variables.
	\item Choose order in which the incomplete variables will be imputed (can influence algorithm convergence) 
		\begin{itemize}
			\item When the missing data pattern is close to monotone: convergence may be quicker if visiting the columns in increasing order of the number of missing data.
		\end{itemize}
	\item Decide setup of the starting imputations and the number of iterations.  
	\item Select the number of imputed datasets.
\end{enumerate}
\vfill
\footnotesize{\url{https://stefvanbuuren.name/fimd/sec-choices.html}}
\end{frame}






\begin{frame}
\frametitle{Reporting}

Guidelines from \footnotesize{\url{https://stefvanbuuren.name/fimd/sec-reporting.html}} (more details and a template of paragraph on the website):
\begin{itemize}
	\item Amount of missing data
	\item Reasons for missingness
	\item Consequences
	\item Method
	\item Software
	\item Number of imputed datasets
	\item Imputation model
	\item Derived variables
	\item Diagnostics
	\item Pooling
	\item Complete-case analysis
	\item Sensitivity analysis
\end{itemize}

\end{frame}


\end{document}


